\chapter{分级阅读路线}

这本书的二十章不必从第一页一路啃到最后一页。数学不是一条单行道，更像一张地铁图：有几条主干线，也有许多换乘站。本附录给你一条“默认路线”和四条“兴趣路线”，你可以按自己的节奏和口味选择，也随时可以下车换乘。

\section{默认路线：四级台阶}

如果你没有特别的偏好，按下面的四级台阶走最稳，每一级都为下一级铺好路。

\textbf{第一级：第 $1$–$4$ 章，打基础。} 从字母、方程和函数出发，把“用符号思考”这件事练熟。这四章是全书的公共语言，后面无论走哪条路线都绕不开，建议按顺序读完，遇到不懂的初等知识随手翻附录 B。

\textbf{第二级：第 $5$–$8$ 章，加深。} 多项式、变换、向量等内容把第一级的工具用得更深。读这一级时你会第一次体会到：同一个对象（比如一个多项式）可以有代数的、几何的好几张面孔，而“换一张面孔看问题”正是数学家最趁手的招式。

\textbf{第三级：第 $9$–$12$ 章，进入高等数学的语言。} 从变化走向无限：极限、导数、积分的直观图景在这里登场。这一级的关键词是“逼近”——允许自己暂时不习惯“无限”的严格说法，先抓住图像和数值例子；严格性会随后跟上，直觉掉队了才真的麻烦。

\textbf{第四级：第 $13$ 章起，按兴趣分流。} 从这里开始，章节之间不再是严格的先后关系：数论、拓扑、几何、概率、数学基础各自成线。读完第三级后，你完全可以“哪里好玩跳哪里”，下面的四条兴趣路线就是为此准备的。

\section{四条兴趣路线}

\textbf{路线一：想解题——第 $1$–$6$、$13$、$14$ 章。} 如果你喜欢“题目刚看时没头绪、解出来一身轻”的快感，这条路线通向数学里最古老的狩猎场：整数。前六章把代数工具磨快，第 $13$、$14$ 章进入数论，那里的问题往往一句话就能听懂（比如关于质数的猜想），却可能难住数学家几百年。建议随身带草稿纸，例题务必亲手算一遍再往下看；解题能力的单位不是“页”，是“草稿纸的张数”。

\textbf{路线二：喜欢图形和空间——第 $3$、$4$、$6$–$12$、$16$、$17$ 章。} 如果你看到图比看到式子更亲切，这条路从坐标与变换出发，经过函数图像、变化率与曲线，最后抵达拓扑和微分几何——那里研究的是“橡皮膜上的几何”和“弯曲空间里的最短路径”。读这条线时多画图，哪怕画得很丑；几何直觉是画出来的，不是看出来的。

\textbf{路线三：喜欢无限和变化——第 $8$–$12$、$15$、$17$ 章。} “无限”是数学里最迷人的悬崖：无限小是什么？无限个正数加起来怎么会有限？这条路线从变化率出发，正面走近极限与无穷级数，再延伸到与之相连的几何与分析话题。走这条线需要一点耐心：每一个关于“无限”的结论，都要追问一句“它是怎么被一步步逼近的”。

\textbf{路线四：喜欢密码、计算和现实应用——第 $2$、$5$、$7$、$13$–$15$、$18$、$19$ 章。} 如果你关心“数学到底有什么用”，这条线从方程与代数结构出发，经过数论（现代密码的地基——你手机里的加密支付就在用质数的性质），再到概率与现实建模。建议配合附录 D 的实验：这条线上的很多想法，用电子表格或几行程序立刻就能玩起来。

\begin{lianjie}
四条路线在第 $13$ 章交汇不是巧合：数论用的工具最“初等”（整除、余数），却连着密码、几何甚至分析。数学各分支之间有许多这样的暗道，本书正文的“连接线”盒子会不断指出它们——发现自己从一条路线溜达到另一条路线上，是读这本书的正常状态，值得高兴，不必纠正。
\end{lianjie}

\section{三条阅读建议}

第一，\textbf{卡住是正常的}。每章只给提示的思考题，卡住几天再回头，往往比当场看答案收获大；数学家们自己也经常在同一个问题上卡几年。第二，\textbf{不必一次读懂}。第一遍可以只追主线故事，把证明留到第二遍；知道“哪里难”本身就是进步。第三，\textbf{动手优先}。遇到“一分钟实验”就停下来做，五分钟的动手常常胜过半小时的瞪眼。

读正文时，另外四个附录是手边工具：附录 A 是证明的语法书，第一次读证明发懵时翻它；附录 B 是初等工具箱，分数、方程、坐标忘了就查；附录 C 是符号字典；附录 D 提供实验器材。它们和本附录一起，构成这本书的“后勤部”——前线在哪一章，后勤随叫随到。

\begin{pandeng}
读完本书想继续攀登的读者，可以按路线找下一座山：解题路线通向竞赛数学与初等数论教材；图形路线通向《什么是数学》和拓扑学普及读物；无限路线通向微积分与数学分析的入门课；应用路线通向密码学和概率论的科普与公开课。关键词在手，图书馆和网络上都是路。
\end{pandeng}

\section*{思考题}

\textbf{观察题}
\begin{sikao}
\item 翻翻目录，挑出你最想读的三章，看看它们各属于哪条路线。提示：一章可以属于多条路线，交叉处往往最好玩。
\end{sikao}

\textbf{证明题}
\begin{sikao}
\item 本附录没有定理要证，但你可以试着向同学论证：为什么第 $1$–$4$ 章被安排在几乎所有路线之前？提示：翻开后面任意一章，数一数它用到了前几章的哪些工具。
\end{sikao}

\textbf{挑战题}
\begin{sikao}
\item 为自己设计一条本附录没有的路线（比如“喜欢数学史和人物故事”），写出停靠章节和每一站值得停的理由。提示：先想终点想看什么风景，再倒推途中需要哪些准备。
\end{sikao}
